% ---------------------------------------------------------------------------
\section{ACE: An Analytical and Differentiable Rate Proxy}
\label{sec:proxy}
% ---------------------------------------------------------------------------

We first define ACE and show that it empirically tracks HEVC bitrate, then provide theoretical guarantees explaining why it does so and why its gradients are reliable.
The definition is closed-form and requires no training data.
The theory establishes three properties---entropy consistency, near-optimality under the standard natural-image model, and gradient alignment with the true HEVC rate---that together justify the use of its gradient as an optimization signal.

\subsection{Definition and Calibration}
\label{sec:definition}

\begin{definition}[ACE]
\label{def:proxy}
Let $\bx \in [0,1]^{H \times W \times C}$ be an image.
Partition $\bx$ into non-overlapping $B \times B$ blocks.
For each block: subtract the block mean; apply the orthonormal 2D DCT-II to obtain $\bC = \bD\tilde{b}\bD^\top$; collect the AC coefficients $\ba = \left(\bC_{ij}\right)_{(i,j)\neq(0,0)}$.
Let $\bc = (c_1,\ldots,c_M)$ be equally-spaced centers on $[-l,\,l]$ and set $\sigma>0$.
The soft-histogram probability for bin $m$ is set as
\begin{equation*}
  \hat{p}_m = \frac{1}{|\ba|}\sum_{a \in \ba}
    \frac{\exp(-\tfrac{(a-c_m)^2}{2\sigma^2})}
         {\sum_{m'}\exp(-\tfrac{(a-c_{m'})^2}{2\sigma^2})},
\end{equation*}
such that the ACE rate of $\bx$ is $\hat{r}(\bx) = -\sum_{m} \hat{p}_m \log_2 \hat{p}_m$.
\end{definition}

ACE replicates HEVC's intra coding pipeline in a closed, differentiable form: block mean subtraction approximates DC intra prediction (HEVC mode 1); orthonormal DCT-II is equivalent to HEVC's transform; AC soft-histogram entropy approximates the CABAC output.
No quantization simulation is included; ACE estimates the entropy of unquantized AC coefficients, which correlates strongly with the quantized coded rate (Table~\ref{tab:calibration}).

ACE is applied independently to each color channel of $\bx$ and the per-channel values are summed to form $\hat{r}(\bx)$.
HEVC codes chroma at reduced resolution with separate QP offsets.
For simplicity, the calibration in Table~\ref{tab:calibration} compares the summed RGB proxy against HEVC's total reported bitrate (luma + chroma combined).
The high within-CRF correlations ($r\ge0.90$) confirm that this simplified color treatment does not impair the gradient signal.

Definition~\ref{def:proxy} requires the specification of different constants; we set $B=8$, $M=64$, $l=1.5$, and $\sigma=3.0/M$.
AC coefficients outside the range $[-l,l]$ are soft-assigned to the nearest boundary bin by the Gaussian kernel.
Empirically, only $0.017\%$ of Kodak~\citep{kodak1993} AC coefficients and $0.030\%$of a 500-image ImageNet~\citep{deng2009imagenet} sample fall outside $[-1.5, 1.5]$, confirming that the range covers the effective support of natural-image distributions.

We validate ACE against the measured HEVC bitrate on all 24 Kodak images at CRF $\in\{22,27,32,37,42,47\}$ (144 (image, CRF) pairs).
ACE is constant per image (QP-agnostic by design), so the 144-pair correlation has two orthogonal sources of variance: the across-image ranking signal and the CRF axis.
Computing Pearson $r$ within each CRF (24 images per CRF) shows $r \in [0.903, 0.981]$ across all six quality levels, with a mean $r = 0.947$.
Similarly, the mean Spearman correlation is $\varrho = 0.932$ (Table~\ref{tab:calibration}).
This demonstrates that gradient descent on $\hat{r}$ moves the image toward a lower HEVC bitrate because images ranked lower by ACE are consistently ranked lower by HEVC at every quality level.

\begin{table}[ht]
  \centering
  \small
  \caption{
  \textbf{ACE's bitrate values correlate in value and ordering with real images.}
  Per-CRF calibration on Kodak (24 images): we compute Pearson $r$ and Spearman $\varrho$ correlations to assess image-level ranking within each quality level.}
  \label{tab:calibration}
  \begin{tabular}{lccccccr}
    \toprule
     & \multicolumn{6}{c}{CRF} & \\
    \cmidrule(lr){2-7}
    Metric & 22 & 27 & 32 & 37 & 42 & 47 & \textbf{Mean} \\
    \midrule
    Pearson $r$      & 0.933 & 0.955 & 0.977 & 0.981 & 0.934 & 0.903 & \textbf{0.947} \\
    Spearman $\varrho$ & 0.911 & 0.953 & 0.984 & 0.949 & 0.919 & 0.875 & \textbf{0.932} \\
    \bottomrule
  \end{tabular}
\end{table}

Gradient steps against ACE are lightweight, requiring approximately $5\,\text{ms}$ on an A100 GPU for a $768\times512$ Kodak image.
Code~\ref{lst:proxy} gives the short, complete PyTorch implementation (using \texttt{torch\_dct.dct\_2d} for the orthonormal DCT-II\footnote{\texttt{pip install torch-dct}; the same result can be obtained via \texttt{torch.fft} without additional dependencies.}).

\begin{tcolorbox}[
    left=0mm, right=0mm, top=0mm, bottom=-1mm,
    title=Code \ref{lst:proxy},
    label={lst:proxy}
]
\begin{lstlisting}[numbers=left, xleftmargin=2em, numbersep=1em, frame=none]
import torch
import torch.nn.functional as F
from torch_dct import dct_2d

def dct_rate_proxy(x, M=64, lam_range=1.5):
    B, C, H, W = x.shape
    blocks = x.unfold(2, 8, 8).unfold(3, 8, 8)
    blocks = blocks.contiguous().view(-1, 8, 8)
    blocks = blocks - blocks.mean((-2, -1), keepdim=True)
    coeffs = dct_2d(blocks, norm='ortho')
    ac = coeffs.view(-1, 64)[:, 1:].reshape(-1)
    sigma = (2 * lam_range) / M
    centers = torch.linspace(-lam_range, lam_range, M, device=x.device)
    log_w = -(ac[:, None] - centers[None, :]).pow(2) / (2 * sigma**2)
    p = F.softmax(log_w, dim=-1).mean(0)
    return -(p * p.clamp(min=1e-10).log2()).sum()
\end{lstlisting}
\end{tcolorbox}

\subsection{Theoretical Analysis}
\label{sec:theory}

Three questions underpin the ACE's practical use.
First, does the soft histogram actually measure entropy, or is it an arbitrary smooth function that merely looks entropy-like?
Second, is the DCT a sensible choice of transform, or would an image-specific basis do substantially better?
Third, and most directly relevant to gradient-based optimization, does $\nabla_{\bx}\hat{r}$ point in the same direction as the gradient of the true HEVC rate?
The three propositions below answer each question in turn.

\paragraph{Does the soft histogram measure entropy?}
The soft-histogram construction in Definition~\ref{def:proxy} is a Gaussian kernel density estimate of the AC coefficient distribution, evaluated on a uniform grid of $M$ centers.
It is not obvious that the resulting quantity relates to Shannon entropy.
The following proposition confirms that it does.

\begin{proposition}
\label{prop:convergence}
Let $p_A$ be bounded and Lipschitz on $[-R,R]$, and let $\sigma>0$ be fixed.
As $M\to\infty$ and $n\to\infty$ with $n\sigma\to\infty$, we have
\[
  \hat{H}_M \;\xrightarrow{P}\; H(A) + \log_2(2\sigma\sqrt{2\pi e}) + o(1).
\]
The bias $B(\sigma)=\log_2(2\sigma\sqrt{2\pi e})$ depends only on $\sigma$, not on $\bx$.  Consequently, (i)~differences $\hat{H}_M(\bx)-\hat{H}_M(\bx^{\prime})$ converge to $H(A_{\bx})-H(A_{\bx^{\prime}})$, and (ii)~$\nabla_{\bx}\hat{r}$ is a consistent estimator of $\nabla_{\bx}H(A)$.
\end{proposition}

\begin{proof}[Proof sketch]
Soft assignment approximates a Gaussian KDE~\citep{parzen1962estimation}.
By standard KDE theory, $\hat{p}_A(a)\to p_A(a)$ uniformly in probability as $n\to\infty$, $\sigma\to0$, $n\sigma\to\infty$ with $\sigma$ fixed.
The Riemann sum over $M$ bins converges to the integral defining $H(A)$, with discretization error $O(1/M)$.
Because $B(\sigma)$ is the same for every image $\bx$, it cancels in differences, and its gradient with respect to $\bx$ is zero.
\end{proof}

In plain terms: the soft histogram is a biased but \emph{consistent} estimator of differential entropy.
The bias is a fixed shift that depends only on the bandwidth $\sigma$, so it plays no role in comparisons between images or in gradient descent.
What the proxy minimizes, in the limit, is precisely the differential entropy of the AC coefficient distribution.

\paragraph{Is the DCT a good transform for this purpose?}
The rate of an entropy coder is minimized by the Karhunen--Lo\`eve transform (KLT), which diagonalizes the signal's covariance and concentrates energy into as few coefficients as possible.
The DCT-II is used in practice because it is data-independent, fast, and---for natural images---a near-perfect approximation to the KLT.
The next proposition quantifies how near.

\begin{proposition}
\label{prop:klt}
Let $\bx$ be a stationary Gaussian AR(1) field.
For $B$-pixel blocks, DCT-II approximates the KLT eigenvalues with error $O(1/B^2)$~\citep{clarke1985transform}.
Because the entropy rate of a Gaussian source is a concave function of log-eigenvalues (via the water-filling formula), this spectral error translates to a rate error $|\hat{r}(\bx) - R^*(\bx)| = O(1/B^2)$.
For $B=8$, $1/B^2 = 1/64 \approx 0.016$, providing a tight approximation.
\end{proposition}

\begin{corollary}
Natural images are well modeled by Markov-1 fields~\citep{habibi1971comparison}.
Minimizing $\hat{r}$ on natural images therefore minimizes a near-optimal proxy for the entropy-coded rate without any codec-specific training.
\end{corollary}

Intuitively, the DCT's near-optimality on natural images is why HEVC uses it: the $O(1/B^2)$ gap is small enough at $B$ equal to $8$ that learning an image-specific basis (as in KLT-based schemes) yields negligible gain.
The same reasoning justifies using the DCT proxy without codec-labeled data.

\paragraph{Does the proxy gradient point the right way?}
Propositions~\ref{prop:convergence} and~\ref{prop:klt} establish that $\hat{r}$ tracks the true entropy-coded rate.
But for gradient-based optimization, the value of $\hat{r}$ matters less than the \emph{direction} of $\nabla_{\bx}\hat{r}$.
A proxy that correlates well in absolute value could still mislead if its gradient points in a different direction from the true rate gradient.

\begin{proposition}[Expected gradient alignment]
\label{prop:grad}
Let $R_\epsilon(\bx)$ denote the mollified HEVC rate, obtained by replacing hard quantization with soft rounding of width $\epsilon$ and CABAC with its differentiable entropy approximation~\citep{balle2018variational}.
$\nabla_{\bx} \hat{r}(\bx)$ shares the block-diagonal DCT-II Jacobian structure of $\nabla_{\bx} R_\epsilon$.
Let $\bx$ be drawn from a distribution over natural images satisfying the Markov-1 model.
Then
\[
  \mathbb{E}_{\bx}\big[\langle \nabla_{\bx}\hat r(\bx),\, \nabla_{\bx} R_\epsilon(\bx)\rangle\big] > 0,
\]
and for each AC frequency coordinate $k$,
\[
  \Pr_{\bx}\!\left(\operatorname{sign}\!\left(\tfrac{\partial \hat r}{\partial a_k}\right) = \operatorname{sign}\!\left(\tfrac{\partial R_\epsilon}{\partial a_k}\right)\right) > \tfrac12.
\]
\end{proposition}

\begin{proof}[Proof sketch]
By the chain rule, $\nabla_{\bx}\hat{r} = (\partial\ba/\partial\bx)^\top\nabla_{\ba}\hat{H}$, where $\partial\ba/\partial\bx$ is block-diagonal with the DCT-II matrix on each block---the same Jacobian structure as $\nabla_{\bx} R_\epsilon$.
It therefore suffices to show $\langle\nabla_{\ba}\hat{H},\nabla_{\ba} H_\epsilon\rangle > 0$ in expectation and coordinatewise sign agreement with probability exceeding one half, where $H_\epsilon$ denotes $R_\epsilon$ viewed as a function of the AC coefficients $\ba$.

Under the Markov-1 model, AC coefficients are approximately independent zero-mean Gaussians with frequency-dependent variances $\sigma_k^2$~\citep{clarke1985transform}.
For a zero-mean Gaussian source, differential entropy is strictly monotone in variance: $\partial H / \partial \sigma_k^2 > 0$.
Increasing $|a_k|$ increases the empirical estimate of $\sigma_k^2$, so $\partial H / \partial |a_k| > 0$.
The same monotonicity holds for $H_\epsilon$ under the soft-rounding entropy model of~\citet{balle2018variational}: the mollified quantization noise adds a fixed variance floor that does not change the sign of the derivative.
Under the idealized independent-Gaussian model, $\nabla_{\ba}\hat{H}$ and $\nabla_{\ba} H_\epsilon$ therefore agree in sign coordinatewise, giving the stated expectation and per-coordinate probability bounds.
Real AC coefficients deviate from this idealization---heavier tails, cross-band correlation, and occasional multimodal blocks at edges and textures---so the coordinatewise agreement is a property that holds on average over natural images rather than at every coordinate of every image, which is why the proposition is stated in expectation and probability rather than pointwise.
\end{proof}

We verify gradient consistency empirically by computing, for each of the 63 AC frequency bands, the centered finite-difference sensitivity of both the proxy and the JPEG rate (Pillow, quality 50) to a $\pm30\%$ scaling of each band's energy, then measuring the cosine similarity of the resulting 63-dimensional gradient vectors.
JPEG uses the same DCT-II structure as the proxy, making it the most direct empirical counterpart.
Across all 24 Kodak images, the mean cosine similarity is $0.48 \pm 0.18$ (all values positive), directly confirming the expectation statement, and the sign of the two gradients agrees in 87\% of AC bands, confirming the per-coordinate probability statement well above the $\tfrac12$ threshold.

This is the central property that makes ACE useful as an optimization signal.
Together, it follows that ACE measures something close to entropy (Proposition~\ref{prop:convergence}), in a way that is a near-optimal rate proxy for the DCT (Proposition~\ref{prop:klt}), and has gradients that push the image in the correct direction to reduce the true coded rate in expectation (Proposition~\ref{prop:grad}).